\chapter{Codici Convoluzionali}
Così come i codici a blocchi, i codici convoluzionali nascono per fronteggiare le condizioni avverse dei canali trasmissivi, introducendo un'opportuna ridondanza in grado di proteggere l'informazione dagli errori.

Tuttavia, mentre i codici a blocchi frammentano il flusso dei dati in segmenti rigidi e indipendenti (ogni parola di $k$ bit genera una codeword di $n$ bit senza alcun legame con i blocchi adiacenti), i codici convoluzionali operano in modo profondamente diverso.

Sebbene siano anch'essi codici \textit{lineari} (strutturati ad albero), il loro codificatore è dotato di \textbf{\textit{memoria}}. È proprio grazie a questa memoria che l'operazione di \textbf{convoluzione} può "sparpagliare" il flusso di dati: la stringa in uscita in un dato istante non dipende solo dai bit attualmente in ingresso, ma è influenzata anche dallo storico dei bit precedenti. Questo meccanismo produce un flusso continuo in cui informazione e bit di controllo risultano indissolubilmente mescolati.
 
\section{Principi e Struttura dei Codici Convoluzionali}

%inserire testo introduttivo qui?

\subsection{L'Operazione di Convoluzione e la Memoria}
Il termine \textit{convoluzionale} deriva dall'omonima operazione matematica ed è il meccanismo centrale tramite il quale questi codici producono l'uscita. Prendendo spunto dall'etimologia della parola, \textbf{avvolgere insieme} è difatti ciò che fanno questi sistemi: aggiungono ridondanza e la intrecciano intimamente con i dati da codificare. Così come in matematica la convoluzione "fonde" due funzioni facendone scorrere una sull'altra, qui il flusso continuo dei dati scorre e si "fonde" con la struttura logica del codificatore.

Da un punto di vista implementativo, se consideriamo $u$ come il flusso di dati in ingresso e $X$ come il flusso codificato in uscita, il processo avviene in modo continuo attraverso questi passaggi:

%TODO: grafico

\begin{itemize}
    \item \textbf{Ingresso:} Dal flusso $u$ vengono letti $k$ bit per volta.
    \item \textbf{Scorrimento (Memoria):} Questi bit entrano in un \textbf{registro a scorrimento} (\textit{shift register}). All'interno di questo registro sono ancora trattenuti i bit elaborati negli istanti precedenti, i quali implementano fisicamente la \textit{memoria} del codificatore.
    \item \textbf{Combinazione:} Una \textbf{funzione combinatoria} (tipicamente una rete di porte logiche XOR) collegata ai vari stadi del registro va a mescolare i nuovi $k$ bit appena entrati con quelli vecchi. Questa operazione produce infine gli $n$ bit in uscita, che andranno a formare il flusso codificato $X$ da trasmettere sul canale.
\end{itemize}

\subsection{Parametri e notazioni}
Elenchiamo parametri e notazioni per definire compiutamente un codificatore \textbf{binario convoluzionale $(n,k)$}
\begin{itemize}

\item \textbf{Data-frame ($k$):} Il gruppo di $k$ bit di informazione che entra nel codificatore in un singolo istante di tempo.

\item \textbf{Code-frame ($n$):} Il gruppo di $n$ bit prodotto in uscita dalla rete logica in ogni istante. Per introdurre ridondanza, si ha sempre $n > k$.

\item \textbf{Frequenza di codifica ($R_c = \frac{k}{n}$):} Detta anche \textit{Code Rate}, indica l'efficienza del codice ed è definita dal rapporto tra i bit di informazione e i bit totali del code-frame

\item \textbf{Lunghezza di vincolo ($N$):} Indica quanti data-frame influenzano la generazione del code-frame. Rappresenta l'ampiezza della "finestra temporale" in cui un bit in ingresso continua a produrre effetti sull'uscita.
\item \textbf{Memoria del codificatore $v=(N-1)k$:} Rappresenta il numero totale di bit \textit{passati}  memorizzati nei registri a scorrimento. Poiché il vincolo totale è $N$, gli stadi puramente di memoria sono $N-1$.
\end{itemize}

\subsection{Codificatore Convoluzionale}
Un codificatore convoluzionale può quindi essere progettato come una \textbf{macchina sequenziale a stati}, con registri a scorrimento con $Nk$ posizioni, raggruppati in $N$ stadi, ciascuno di $k$ elementi. A loro volta i registri sono collegati a un circuito combinatorio collegato al registro di uscita; insieme al vincolo, il modo in cui è progettato il circuito determina il comportamento del codificatore.

Per determinare il funzionamento del codificatore, faremo uso dei \textit{generatori}.
%TODO: grafico

\subsection{Sequenze Generatrici}
Mentre nei codici lineari a blocchi, la codifica è determinata da matrici generatrici o da polinomi generatori, nei codici convoluzionali la codifica è specificata dalle \textbf{sequenze generatrici} (\textit{generatori}).

Per determinare il comportamento di un codice convoluzionale $(n,k,N)$ ci si affida a $nk$ \textit{vettori di bit}, di lunghezza $N$, che identificano matematicamente (e \textit{fisicamente} identificando quali commutatori attivare) il codificatore. Definiscono in modo univoco il legame tra il flusso di dati in ingresso e i bit prodotti in uscita.

Sia $[m_1,m_2,m_3,\dots]$ il flusso di messaggi in ingresso al codificatore.

I generatori relativi all'uscita i-esima del codificatore sono indicati con:
\[ g_{ij}=[g_{ij}^1,g_{ij}^2,\dots,g_{ij}^N] \quad i=1,\dots,n \quad j=1,\dots,k\]

Sia $u_j^{(t)}$ il j-esimo ingresso al tempo $t$, con $j = 1,\dots , k$

Sia $x_i^{(t)}$ la i-esima uscita al tempo $t$, con $i = 1,\dots , n$
\[x_i^{(t)}= \sum_{j=1}^k\sum_{l=1}^N g_{ij}^lu_j^{(t-l)}\quad i=1,\dots,n\]


Per un codificatore con vincolo di dimensione $N$, i generatori sono definiti dalle $n$ definizioni:
\[g_i= [g_i^1,g_i^2,\dots,g_i^N]\quad i=1,\dots,n\]

Il valore del coefficiente $g_i^l$ è pari a $1$ quando lo stadio $l$-esimo del registro in ingresso è \textbf{\textit{connesso}} al sommatore, la cui uscita è a sua volta connessa allo stadio $i$-esimo del registro in uscita. La sequenza in uscita all'$i$-esimo sommatore (assumendo $k=1$ singolo bit in ingresso), è quindi data dalla convoluzione:
\[x_i^{(t)}= \sum_{l=1}^N g_i^{l}u^{(t-l)}\quad i=1,\dots, n\]

\subsection{Rappresentazioni Grafiche}
Si possono usare diverse rappresentazioni grafiche per rappresentare un codificatore convoluzionale:
\begin{itemize}
    \item rappresentazione \textit{ad albero}
    \item rappresentazione \textit{a traliccio}
    \item diagramma degli stati 
\end{itemize}
\subsubsection{Diagramma a Stati}
Il diagramma degli stati è un \textbf{grafo} composto da $2^v = 2^{k(N-1)}$ nodi, ognuno dei quali è uno dei possibili stati in cui può trovarsi il codificatore.

\noindent Con il grafo si può modellare completamente il comportamento di un codice convoluzionale.

\noindent Ogni stato è identificato dagli $i$ bit nei registri a scorrimento, e dai $k$ bit in ingresso, che di concerto determinano gli $n$ bit in uscita.

\noindent I nodi sono connessi tra loro tramite archi, ognuno dei quali identifica un data-frame in ingresso. Per data-frame di dimensione $k$, si avranno $2^k$ rami uscenti per ogni nodo

\noindent La transizione da uno stato ad un altro è associata alla emissione di un code-frame di $n$ bit.

\subsubsection{Diagramma a Traliccio (Trellis)}

Il diagramma \textbf{a traliccio} può essere considerato come un'estensione \textit{temporale} di un diagramma a stati; aggiunge infatti un indice temporale $i$. È uno strumento utile per analizzare la decodifica.

\noindent Per ogni indice temporale, il numero degli stati è determinato dalla memoria $v$ del codificatore; gli stati saranno quindi $2^v = 2^{k(N-1)}$.

\noindent Nel passaggio da un indice temporale al successivo, gli stati sono connessi tramite gli archi delle possibili transizioni di stato, ognuna delle quali è determinata da una possibile sequenza in ingresso.

\begin{esempio}{Codice Convoluzionale $(2,1,3)$} 
Consideriamo un codificatore caratterizzato da $n=2$, $k=1$, $N=3$. 
I due generatori associati alle uscite sono $g_1 = [110]$ e $g_2 = [011]$. 
Le equazioni per i due bit di uscita all'istante $i$ risultano:
\begin{align*}
    x_1 &= u_i + u_{i-1} \\
    x_2 &= u_{i-1} + u_{i-2}
\end{align*}
Lo stato del sistema è determinato dagli ultimi due bit in ingresso memorizzati nei registri ($u_{i-1}, u_{i-2}$). Al variare dell'ingresso $u_i$, la macchina transita tra i 4 stati possibili ($00, 10, 01, 11$), emettendo i relativi bit di uscita $x_1 x_2$ e tracciando un percorso specifico sul diagramma a stati e sul traliccio.
\end{esempio}
\subsection{Capacità Correttiva: Distanze del Codice}

Le capacità di rivelazione e correzione degli errori di un codice convoluzionale sono vincolate alla stessa metrica fondamentale utilizzata per i codici a blocchi: la \textbf{distanza di Hamming}.

Tuttavia, emerge una differenza sostanziale nel modo in cui questa metrica viene valutata. Non elaborando \textit{singole codeword} isolate, ma operando su \textit{flussi continui}, è necessario considerare la distanza tra le \textbf{intere sequenze} codificate $X$, generate a partire dai diversi flussi di informazione in ingresso.

Poiché i codici convoluzionali mantengono la proprietà di \textbf{linearità,} calcolare la distanza minima tra due sequenze valide qualsiasi equivale semplicemente a valutare il \textbf{peso di Hamming} dell'intera sequenza (ovvero la sua distanza rispetto alla sequenza nulla composta da soli zeri).
Non cambia quindi il paradigma che guida la decodifica, il criterio di \textbf{Massima Verosimiglianza} (\textit{Maximum Likelihood}).

Si introducono due metriche specifiche per i codici convoluzionali, la distanza colonna e la distanza libera.

\subsubsection{Distanza Colonna}
Consideriamo due sequenze codificate fino all'istante $i$ in cui divergono alla prima diramazione del traliccio.

La distanza colonna $d_c(i)$ è la misura della minima distanza di Hamming tra \textbf{tutte le coppie} di sequenze di questo tipo.

Per calcolare $d_c(i)$ si considerano le distanze rispetto alla \textbf{sequenza di tutti zeri}. In questo modo, sfruttando la linearità del codice, la distanza colonna coincide con il \textbf{peso minimo} delle sequenze di codice che divergono alla prima ramificazione dalla \textit{sequenza zero}.

La $d_c(i)$ è una funzione \textbf{non decrescente} della profondità $i$, ovvero:
$$ d_c(i) \le d_c(i+1)$$

L'apparente limitazione di valutare la distanza rispetto alla sola \textbf{sequenza nulla} (tutti zeri) è giustificata dalla \textbf{linearità} del codice. Poiché la distanza di Hamming tra due sequenze codificate $A$ e $B$ è identica al peso di Hamming della loro somma modulo-2 ($A \oplus B$), e dato che in un codice lineare tale somma è anch'essa una sequenza valida, la sequenza di tutti zeri funge da \textbf{riferimento universale}. 

Trovare il percorso che diverge e accumula il minor peso rispetto allo zero equivale a trovare la distanza minima tra \textbf{qualsiasi} coppia di sequenze del codice, senza dover confrontare ogni possibile combinazione di messaggi.

\subsubsection{Distanza Libera}
Estendendo l'analisi della distanza colonna verso una profondità infinita, definiamo la \textbf{distanza libera} come il valore limite:
\[ d_l = \lim_{i \to \infty} d_c(i) \]

Rappresenta la minima distanza di Hamming tra due sequenze codificate di lunghezza arbitraria (\textit{infinitamente lunghe}). Grazie alla linearità, $d_l$ coincide con il peso minimo di un percorso che:
\begin{enumerate}
    \item Diverge dal percorso nullo (stato tutto zero) al primo istante.
    \item Attraversa il traliccio \textit{accumulando peso}.
    \item Si ricongiunge \textit{alla fine} allo stato nullo iniziale.
\end{enumerate}


Il calcolo della distanza libera può quindi essere effettuato senza dover considerare il peso $w$ di tutte le possibili sequenze $w(X)$:

Si identifica sul traliccio il percorso a \textbf{peso minimo} che si \textit{separa} dal percorso \textit{nullo}, per poi ricongiungersi, \textit{stabilmente} con quest'ultimo.

Come già detto, $d_l$ sarà il \textbf{minimo peso} di questo insieme di sequenze.

Operativamente si suppone di terminare una sequenza con una serie di bit $0$ in modo da riportare il registro a scorrimento al suo stato iniziale, composto di tutti $0$. In questo modo si eliminano alcuni rami dalle $L$ transizioni finali nel traliccio del codice.

In fase di progettazione, fissati il rapporto di codifica $R_c = k/n$ e la lunghezza del vincolo $N$, il codice \textbf{ottimo} è quello che garantisce la \textbf{massima distanza libera} possibile, minimizzando così la probabilità di errore residua dopo la decodifica.

%TODO: esempio

\subsection{Funzione Generatrice}
Mentre la distanza libera indica il \textit{caso peggiore}, si introduce la \textbf{funzione generatrice} per fornire l'intera distribuzione dei pesi del codice. Essa può essere vista come la funzione \textit{di trasferimento} del sistema, valutata rispetto alle transizioni tra gli stati del codificatore.

La funzione generatrice fornisce indicazioni fondamentali sia sulla distanza libera, sia sulle probabilità di errore che si possono incontrare in fase di decodifica.

\section{Decodifica dei Codici Convoluzionali} 
Sono tre i metodi di decodifica principali per i codici convoluzionali:
\begin{itemize}
    \item \textbf{Algoritmo di Viterbi:} implementa una decodifica ottima a massima verosimiglianza, richiede hardware \textit{complesso};

    \item \textbf{Decodifica Sequenziale:} realizza un compromesso tra complessità implementativa e prestazioni della decodifica;

    \item \textbf{Decodifica a Retroazione:} decodifica che non richiede hardware \textit{complesso}, ma offre prestazioni inferiori;
\end{itemize}

\subsection{Metriche per la Decodifica a Massima Verosimiglianza su Traliccio}
Similmente ai codici lineari a blocchi, anche per i codici convoluzionali il criterio di \textbf{massima verosimiglianza} viene espletato selezionando di volta in volta la sequenza di codice a \textbf{distanza minima}.

Essendo la \textit{possibile sequenza codificata} un cammino sul traliccio del codice, il decodificatore individua il cammino a distanza minima dalla sequenza ricevuta.

Si possono utilizzare due metriche:
\begin{itemize}
    \item \textbf{Distanza di Hamming:} nel caso di decodifica \textit{hard};
    \item \textbf{Distanza Euclidea:} nel caso di decodifica \textit{soft};
\end{itemize}

\subsection{L'Algoritmo di Viterbi}
In \textit{linea teorica}, un decodificatore a massima verosimiglianza dovrebbe analizzare l'intera sequenza ricevuta $Y$, trovando il percorso \textit{ammissibile} a minore distanza di Hamming da $Y$, ma in casi reali non è una strategia percorribile, il costo sarebbe proibitivo perché il numero di percorsi possibili cresce esponenzialmente con la lunghezza della sequenza.

L'algoritmo proposto da Andrew Viterbi nel 1967, rende possibile la ricerca del cammino ottimale senza ricorrere a una ricerca esaustiva.

Questo risultato si ottiene attraverso due pilastri fondamentali:
\begin{itemize}
    \item L'introduzione della \textbf{metrica di ramo} (\textit{branch metric}), che permette di valutare la distanza locale ad ogni transizione;
    \item La selezione sistematica dei \textbf{percorsi sopravvissuti}, che consiste nel \textit{potare} precocemente i cammini che non possono più ambire ad essere ottimali, evitando così di esplorare rami inutili.
\end{itemize}



\subsubsection{Metrica di Ramo e Metrica di Percorso}

Consideriamo il caso in cui in ingresso al registro a scorrimento arriva un solo bit ($k=1$).

Ogni nodo del traliccio al generico stato $j$, all'istante $i$,è indicato come:
\[S_i^j \quad j=1,2,\dots,2^{N-1}\]

\textbf{Metrica del ramo:}
\[\lambda(S_i^j,S_{i+1}^l)\]
La distanza (di Hamming o euclidea) tra gli $n$ bit \textbf{ricevuti} al tempo $i$-esimo e gli $n$ bit \textbf{emessi} con la transizione dallo stato $j$-esimo allo stato $l$-esimo:
\[S_i^j \Rightarrow S_{i+1}^l\]
La metrica del ramo è ovviamente \textbf{additiva}, infatti ogni percorso $P$ composto da $M$ transizioni avrà la \textbf{metrica di percorso} (lunghezza del cammino) definita come \textbf{somma delle metriche di ramo calcolate per ogni transizione di stato che lo compone}:

\[\Gamma(P)= \sum_{k=0}^{M-1}\lambda \left(S_{i+k}^{j_k}, S_{i+k+1}^{j_{k+1}}\right)\]

\subsubsection{Concetto di Percorsi Sopravvissuti}
Serve un ultimo strumento per l'algoritmo di Viterbi utile per il processo di \textit{scrematura} dei percorsi implicito nell'algoritmo stesso.
Osservando il traliccio si nota che in ogni istante temporale, in un singolo nodo (stato) convergono sempre $2^k$ rami provenienti da stati precedenti.
Poiché questi percorsi, da questo specifico stato in poi, dovranno necessariamente condividere lo stesso identico futuro, solo il cammino che ha accumulato finora la \textbf{metrica di percorso minore} (distanza minima) può \textit{matematicamente} fare parte del cammino ottimo globale.

Il \textbf{percorso sopravvivente} $\hat{P}(S_i^j)$ sarà, tra quelli che giungono in quel nodo, quello  con la minore lunghezza del cammino.

\subsubsection{Procedura Operativa dell'Algoritmo}
\begin{enumerate}
    \item \textbf{Inizializzazione:} $i = 0$ e $\Gamma(S_i^j), \space j=1,2,\dots,2^{N-1}$; si assume che lo stato iniziale e quello finale dopo $M$ passi siano composti da tutti \textit{zeri} (reset dei registri).

    \item \textbf{Calcolo metriche di ramo:} Per $i=i+1$ si calcolano le $2^{N-1}$ metriche di ramo $\lambda(S_{i-1}^j,S_i^l)$ per ognuno dei $2^{N-1}$ stati.

    \item \textbf{Selezione dei percorsi sopravviventi}: per ogni stato $S_i^j$ si calcolano le lunghezze dei percorsi incidenti e si memorizza il percorso sopravvivente $\hat{P}(S_i^j)$ e la nuova metrica di stato:
    $$\Gamma(S_i^j) = \Gamma(\hat{P}(S_i^j))$$

    \item \textbf{Iterazione:} si ripetono i passi precedenti fino al passo $M$.

    \item \textbf{Decodifica:} si decodifica la sequenza ricevuta come la sequenza associata alle transizioni di stato dell'unico percorso sopravvivente.
\end{enumerate}

\subsection{Limiti e Complessità Computazionale}
Nonostante l'algoritmo di Viterbi offra una decodifica ottima a massima verosimiglianza, la sua applicazione pratica è fortemente limitata dalla complessità computazionale e dai requisiti di memoria, che crescono esponenzialmente al crescere dei parametri del codificatore. Solamente di recente sono state realizzate implementazioni hardware ad-hoc.

Per un codice convoluzionale $(n, k, N)$, ad ogni passo (istante temporale) sul traliccio:
\begin{itemize}
    \item Si hanno $2^{k(N-1)}$ stati possibili. Il decodificatore deve mantenere attivi e memorizzare $2^{k(N-1)}$ \textbf{cammini sopravviventi} (con relative sequenze di bit) e i corrispondenti $2^{k(N-1)}$ \textbf{valori di metrica di percorso}.
    \item Per ogni stato si hanno $2^k$ rami provenienti dall'istante precedente. Il decodificatore deve quindi calcolare $2^k$ \textbf{metriche di ramo} per ogni singolo stato, aggiornare le metriche di percorso e confrontarle per selezionare il percorso sopravvivente.
\end{itemize}

A causa di questa dipendenza esponenziale sia dal numero di bit in ingresso $k$, sia dalla lunghezza del vincolo $N$, l'implementazione hardware dell'algoritmo di Viterbi diventa rapidamente proibitiva. Per questo motivo, nelle applicazioni reali (come le comunicazioni spaziali o satellitari), l'utilizzo di Viterbi è limitato a codici con valori di $k$ e $N$ dell'ordine di grandezza di $10$. Per codici più complessi si ricorre solitamente a metodi sub-ottimi, come la decodifica sequenziale.